\documentclass[10pt,a4paper]{article}
\usepackage[margin=2cm]{geometry}
\usepackage[T1]{fontenc}
\usepackage[utf8]{inputenc}
\usepackage[spanish]{babel}
\usepackage{array}
\usepackage{enumitem}

\setlength{\parindent}{0pt}
\setlist[itemize]{leftmargin=1.2em,itemsep=0.2em}

\title{Informe Técnico: Hoja de Cálculo con Matriz Dispersa}
\author{Leonardo Gabriel Sanchez Terrazos\\202410273}
\date{}

\begin{document}
\maketitle

\section*{Justificación de la Estructura de Datos}
La aplicación modela la hoja de cálculo como una matriz dispersa con listas enlazadas cruzadas. Esta representación evita almacenar celdas vacías y permite recorrer solo los elementos realmente presentes. En este esquema, insertar o eliminar un nodo tiene costo $O(k)$, donde $k$ es la cantidad de elementos de la fila o columna involucrada, lo que resulta más eficiente que recorrer una matriz densa completa cuando la mayoría de celdas están vacías.

\section*{Análisis de Complejidad}
La siguiente estructura queda preparada para documentar las complejidades temporales y espaciales de las 13 operaciones del sistema.

\vspace{0.5em}
\begin{tabular}{|>{\raggedright\arraybackslash}p{0.56\linewidth}|>{\centering\arraybackslash}p{0.18\linewidth}|>{\centering\arraybackslash}p{0.18\linewidth}|}
\hline
\textbf{Operación} & \textbf{Tiempo} & \textbf{Espacio} \\ \hline
Inserción de celda & --- & --- \\ \hline
Eliminación de celda & --- & --- \\ \hline
Actualización de celda & --- & --- \\ \hline
Búsqueda de celda & --- & --- \\ \hline
Recorrido de fila & --- & --- \\ \hline
Recorrido de columna & --- & --- \\ \hline
Creación de matriz & --- & --- \\ \hline
Limpieza de matriz & --- & --- \\ \hline
Carga de archivo & --- & --- \\ \hline
Guardado de archivo & --- & --- \\ \hline
Visualización de hoja & --- & --- \\ \hline
Redimensionamiento lógico & --- & --- \\ \hline
Validaci\'on de coordenadas & --- & --- \\ \hline
\end{tabular}

\end{document}